%% kapa: exam=1 round=21 year=2010 period=2 form=A subject=accounting subject_ko=회계학 seq=7 number=7
%% source: Q-Net 감정평가사 기출문제 (cst00309) artlSeq=5024181 file=extracted/1008/5024181/2010년 제21회 감정평가사 1차시험 기출문제/2교시.hwp
%% conversion: kordoc 4.24.51 -> tools/gen1.py (auto); review=proofread-gpt-sol; flags=none
\begin{question}{7}
한국채택국제회계기준에 의하면 재무상태표에 구분 표시해야 할 항목에 ‘투자부동산’이 포함된다. 다음 중 투자부동산의 예로 옳지 않은 것은?
\begin{choices}
  \item 제3자를 위하여 건설 또는 개발 중인 부동산
  \item 직접 소유하고 운용리스로 제공하고 있는 건물
  \item 운용리스로 제공하기 위하여 보유하고 있는 미사용 건물
  \item 장래 사용목적을 결정하지 못한 채로 보유하고 있는 토지
  \item 미래에 투자부동산으로 사용하기 위하여 건설 또는 개발 중인 부동산
\end{choices}
\end{question}
\answer{1}
